\chapter*{Posicionamiento bibliográfico mínimo}
\addcontentsline{toc}{chapter}{Posicionamiento bibliográfico mínimo}

\section*{Objeto}

\begin{observation}
Esta nota fija la separación mínima entre antecedentes externos, construcciones propias de MADELEINE y relaciones de comparación utilizadas en la presente fase de publicación. No constituye una revisión exhaustiva del estado del arte ni pretende reconstruir retrospectivamente el origen del núcleo MADELEINE a partir de literatura externa.
\end{observation}

\section*{Antecedente externo actualmente incorporado}

\begin{observation}
La noción estándar de cadena de Markov en tiempo discreto se toma aquí como referencia externa de terminología y comparación \parencite{norris1997markov}. Esta referencia documenta únicamente el marco externo correspondiente y no demuestra ningún resultado propio de MADELEINE.
\end{observation}

\section*{Construcción MADELEINE}

\begin{observation}
LAB-FEN-001, sus controles, sus métricas, las ablaciones, el atlas paramétrico, DMO-004, HISTORY-DEPTH-v1 y HISTORY-DEPTH-ATLAS-v1 pertenecen a la construcción metodológica interna de MADELEINE en esta fase. Su soporte efectivo permanece en las especificaciones, ejecutables, resultados estructurados, actas, hashes e historial Git del proyecto.
\end{observation}

\section*{Relación entre antecedente y construcción}

\begin{observation}
Los sustitutos Markov-1 se utilizan en LAB-FEN-001 como referencias nulas operacionales dentro de métricas concretas. Esta utilización no identifica la profundidad operativa \(k^\ast\) con un orden Markoviano universal ni convierte HISTORY-DEPTH-v1 en una reformulación de la teoría estándar de cadenas de Markov.
\end{observation}

\begin{observation}
La relación bibliográfica declarada es, por tanto, de referencia terminológica y comparación metodológica restringida. No se declara dependencia fundacional del núcleo MADELEINE respecto de la teoría de cadenas de Markov.
\end{observation}

\section*{Límites de posicionamiento}

\begin{itemize}
\item no se presenta la bibliografía externa como validación de MADELEINE;
\item no se atribuyen a fuentes externas los resultados numéricos de LAB-FEN-001;
\item no se identifica semejanza terminológica con equivalencia formal;
\item no se declara que MADELEINE sustituya teorías clásicas;
\item no se formula todavía un estado del arte exhaustivo;
\item no se incorporan referencias sin una función pública concreta y verificable.
\end{itemize}

\section*{Criterio de ampliación futura}

\begin{observation}
Nuevas fuentes sólo deberán incorporarse cuando exista una oración pública concreta cuya comprensión, procedencia o comparación disciplinada requiera esa referencia. La ampliación bibliográfica no debe utilizarse para inflar artificialmente el corpus ni para crear dependencias retrospectivas inexistentes.
\end{observation}
